\documentclass{article} % say
\usepackage{amsmath}
\usepackage{amssymb}
\usepackage{gensymb}
\usepackage{hyperref}
\usepackage{tikz}
\usepackage{tkz-euclide}
\usepackage{listings}
\usetikzlibrary {babel,angles,quotes,positioning}
\usepackage[utf8]{inputenc}
\usepackage[russian]{babel}

\usepackage{mathtools}
\DeclarePairedDelimiter\ceil{\lceil}{\rceil}
\DeclarePairedDelimiter\floor{\lfloor}{\rfloor}
\begin{document}

\textbf{Задача}: Даны два натуральных числа $n$ и $k$, при чём $k \leqslant n \leqslant 10^6$. Необходимо составить массив длинной не больше $25$, чтоб из его элементов можно было собрать любую сумму от $1$ до $n$ кроме $k$.

\clearpage

\textbf{Решение}: 

Для начала составим массив, из которого можно собрать все суммы от $1$ до $k-1$ и только их. Пусть мы получили массив, из которого можно собрать все суммы от $1$ до $S$ и только их. Если мы добавим в этот массив число $S' \leqslant S + 1$, то легко видеть, что в полученном массиве достижимы суммы от $1$ до $S + S'$ и только они. Таким образом, мы можем добавлять степени двойки, пока очередная степень не даст суммe, большую либо равную $k$. Тогда вместо этой степени добавим число $k - 1 - S$, где $S$~--- сумма элементов массива на соответсвующем шаге.

Теперь нужно расширить массив, чтоб получить все суммы от $k+1$ до $n$. Будем добавлять только числа, большие $k$. Тем самым мы получим, что сумму $k$ собрать нельзя (ведь все числа полученного массива, не превосходящие $k$, дают сумму $k-1$ по построению). Для начала добавим в массив число $k+1$. После этого нам достижимы все суммы от $1$ до $2k$, исключая $k$. Затем добавим число $2k+1$: теперь нам доступны все суммы от $1$ до $4k + 1$, исключая $k$ и $3k+1=k + (2k + 1)$. Снова добавим в массив число, равное сумме элементов массива плюс один, пока не доберёмся до суммы, не меньшей $n$. Обозначим добавленные числа как $a_1, \ldots, a_m$. В полученном массиве нельзя получить суммы вида $k + a_i$ (для $i > 1$). Легко видеть, что добавив в массив число $k+1$, мы решаем озвученную проблему: ведь сумма $a_i + k - (k + 1) = a_i - 1$ достижима.

\lstinputlisting[language=C++]{main.cpp}

\textbf{Источник}: \url{https://codeforces.com/contest/1965/problem/B}

\end{document}


